%!TEX TS-program = xelatex
\documentclass[11pt]{article}
\usepackage[margin=1in]{geometry}
\usepackage{fontspec}

\newfontfamily\kurdishfont[Script=Arabic]{Amiri}
% For code listings: FreeMono keeps brackets clear and supports Kurdish text.
\newfontfamily\codemonofont{FreeMono}
\newcommand{\codeleftparen}{{\codemonofont(}}
\newcommand{\coderightparen}{{\codemonofont)}}
\newcommand{\codeleftbracket}{{\codemonofont[}}
\newcommand{\coderightbracket}{{\codemonofont]}}
\newcommand{\codeleftbrace}{{\codemonofont\char`\{}}
\newcommand{\coderightbrace}{{\codemonofont\char`\}}}

\usepackage{hyperref}
\usepackage{booktabs}

% Core packages
\usepackage{amsmath,amssymb}
\usepackage{tikz-cd}
\usepackage{multicol}
\usepackage{xcolor}
\usepackage{listings}

% Code listings
\lstdefinestyle{pythoncode}{
  language=Python,
  basicstyle=\footnotesize\codemonofont,
  keywordstyle=\color{blue!70!black},
  commentstyle=\color{green!40!black},
  stringstyle=\color{orange!50!black},
  showstringspaces=false,
  breaklines=true,
  breakatwhitespace=false,
  columns=fullflexible,
  keepspaces=true,
  tabsize=4,
  frame=single,
  numbers=left,
  numberstyle=\scriptsize\codemonofont\color{gray},
  stepnumber=1,
  numbersep=10pt,
  extendedchars=false,
  literate=
    {(}{{\codeleftparen}}1
    {)}{{\coderightparen}}1
    {[}{{\codeleftbracket}}1
    {]}{{\coderightbracket}}1
    {\{}{{\codeleftbrace}}1
    {\}}{{\coderightbrace}}1
}

\lstdefinestyle{bashcode}{
  language=bash,
  basicstyle=\footnotesize\codemonofont,
  keywordstyle=\color{blue!70!black},
  commentstyle=\color{green!40!black},
  stringstyle=\color{orange!50!black},
  showstringspaces=false,
  breaklines=true,
  breakatwhitespace=false,
  columns=fullflexible,
  keepspaces=true,
  tabsize=4,
  frame=single,
  numbers=left,
  numberstyle=\scriptsize\codemonofont\color{gray},
  stepnumber=1,
  numbersep=10pt,
  extendedchars=false,
  literate=
    {(}{{\codeleftparen}}1
    {)}{{\coderightparen}}1
    {[}{{\codeleftbracket}}1
    {]}{{\coderightbracket}}1
    {\{}{{\codeleftbrace}}1
    {\}}{{\coderightbrace}}1
}

% Paragraphs
\setlength{\parindent}{0pt}
\setlength{\parskip}{1\baselineskip}

\title{BERT Vector Generation and Visualization Toolkit}
\author{Data Engineering Team}
\date{\today}

\begin{document}
\maketitle

\section*{Overview}

This document contains a comprehensive toolkit of scripts utilized for generating vector embeddings from the Kurdish medical dataset using the BERT model, and visualizing the semantic relationships between these embeddings. Below is a summary table containing links to the code for each step in the pipeline, alongside an explanation of its purpose.

\begin{table}[htbp]
\centering
\small
\renewcommand{\arraystretch}{1.3}
\begin{tabular}{@{} l l p{9cm} @{}}
\toprule
\textbf{Script Name} & \textbf{Section Link} & \textbf{Purpose / Why We Use It} \\
\midrule

\multicolumn{3}{@{}l}{\textbf{Vector Pipeline}} \\
\midrule
Vector Generator & \hyperref[sec:generator]{Section \ref*{sec:generator}} & Generates vector embeddings for a batch of texts using the BERT large model. \\
Vector Visualization & \hyperref[sec:visualize]{Section \ref*{sec:visualize}} & Computes cosine similarity between vectors and builds an interactive pyvis network graph. \\
\addlinespace

\multicolumn{3}{@{}l}{\textbf{Setup \& Environment}} \\
\midrule
Setup Script & \hyperref[sec:setup]{Section \ref*{sec:setup}} & Prepares the environment by installing dependencies and caching the BERT model. \\
Model Loader & \hyperref[sec:loader]{Section \ref*{sec:loader}} & A dedicated script to trigger the Hugging Face download of the BERT model for caching. \\
Requirements & \hyperref[sec:requirements]{Section \ref*{sec:requirements}} & Lists the necessary Python packages for the pipeline. \\
Environment Variables & \hyperref[sec:env]{Section \ref*{sec:env}} & Configuration file for specifying model and paths. \\
\bottomrule
\end{tabular}
\caption{Vector generation and setup scripts categorized by their functional domain.}
\end{table}

\clearpage

\section{Vector Generator} \label{sec:generator}
\textbf{Purpose:} We use this script to generate vector embeddings for the Kurdish medical dataset using the \texttt{bert-large-cased} model. It includes batching, checkpointing, and GPU acceleration.

\begin{lstlisting}[style=pythoncode, caption={vector\_generator.py}]
import json
import os
import logging
from typing import List, Dict
from tqdm import tqdm
import torch
from transformers import AutoTokenizer, AutoModel

# Setup Logging
logging.basicConfig(
    level=logging.INFO,
    format='%(asctime)s - %(levelname)s - %(message)s',
    handlers=[
        logging.FileHandler("vector_pipeline.log"),
        logging.StreamHandler()
    ]
)

# Configuration
# We use bert-large-cased as requested for BERT
MODEL_NAME = "bert-large-cased"
DATASET_PATH = "../just_data_vector/kurdish_medical_corpus_kmc.json"
OUTPUT_PATH = "kurdish_medical_vectors.jsonl"
CHECKPOINT_PATH = "processed_ids.txt"

class VectorGenerator:
    def __init__(self):
        logging.info(f"Downloading/Loading Model: {MODEL_NAME}")
        # Initialize tokenizer and model
        self.tokenizer = AutoTokenizer.from_pretrained(MODEL_NAME)
        self.model = AutoModel.from_pretrained(MODEL_NAME)
        
        # Check if GPU is available
        self.device = torch.device("cuda" if torch.cuda.is_available() else "cpu")
        self.model.to(self.device)
        self.model.eval()
        logging.info(f"Model loaded and moved to {self.device}")
        
        self.processed_ids = self._load_checkpoints()

    def _load_checkpoints(self) -> set:
        if os.path.exists(CHECKPOINT_PATH):
            with open(CHECKPOINT_PATH, 'r') as f:
                return set(line.strip() for line in f)
        return set()

    def _save_checkpoint(self, entry_id: str):
        with open(CHECKPOINT_PATH, 'a') as f:
            f.write(f"{entry_id}\n")
        self.processed_ids.add(entry_id)

    def clean_text(self, text: str) -> str:
        if not text:
            return ""
        # Basic cleanup
        text = text.replace('\n', ' ')
        text = ' '.join(text.split())
        return text

    def generate_vectors_batch(self, texts: List[str]) -> List[List[float]]:
        """
        Generates vector embeddings for a batch of texts using XLM-R.
        """
        if not texts:
            return []
            
        try:
            # Tokenize input text in batches
            inputs = self.tokenizer(
                texts,
                return_tensors="pt",
                truncation=True,
                max_length=512,
                padding=True
            ).to(self.device)
            
            # Generate embeddings
            with torch.no_grad():
                outputs = self.model(**inputs)
                
            # Mean pooling of the last hidden states
            attention_mask = inputs['attention_mask'].unsqueeze(-1).expand(outputs.last_hidden_state.size()).float()
            sum_embeddings = torch.sum(outputs.last_hidden_state * attention_mask, 1)
            sum_mask = torch.clamp(attention_mask.sum(1), min=1e-9)
            mean_pooled = sum_embeddings / sum_mask
            
            # Convert to list of lists of floats
            vectors = mean_pooled.cpu().numpy().tolist()
            return vectors
            
        except Exception as e:
            logging.warning(f"Vector generation failed for batch... Error: {e}")
            return [[] for _ in texts]

    def store_vectors_batch(self, entry_ids: List[str], vectors: List[List[float]]):
        """
        Appends a batch of generated vectors to the JSONL file.
        """
        with open(OUTPUT_PATH, 'a', encoding='utf-8') as f:
            for entry_id, vector in zip(entry_ids, vectors):
                if not vector:
                    continue
                record = {
                    "id": entry_id,
                    "vector": vector
                }
                f.write(json.dumps(record, ensure_ascii=False) + '\n')

    def process_dataset(self, batch_size: int = 32):
        logging.info(f"Loading dataset from {DATASET_PATH}...")
        try:
            with open(DATASET_PATH, 'r', encoding='utf-8') as f:
                data = json.load(f)
        except Exception as e:
            logging.error(f"Could not load JSON: {e}")
            return

        logging.info(f"Processing {len(data)} entries with batch_size={batch_size}. Found {len(self.processed_ids)} already processed.")
        
        batch_ids = []
        batch_texts = []
        
        for entry in tqdm(data):
            entry_id = entry.get('id')
            if not entry_id or entry_id in self.processed_ids:
                continue
                
            raw_text = entry.get('response', '')
            if not raw_text:
                self._save_checkpoint(entry_id)
                continue
                
            cleaned_text = self.clean_text(raw_text)
            
            batch_ids.append(entry_id)
            batch_texts.append(cleaned_text)
            
            if len(batch_ids) >= batch_size:
                vectors = self.generate_vectors_batch(batch_texts)
                self.store_vectors_batch(batch_ids, vectors)
                
                # Save checkpoints
                for b_id in batch_ids:
                    self._save_checkpoint(b_id)
                    
                batch_ids = []
                batch_texts = []
                
        # Process remaining
        if batch_ids:
            vectors = self.generate_vectors_batch(batch_texts)
            self.store_vectors_batch(batch_ids, vectors)
            for b_id in batch_ids:
                self._save_checkpoint(b_id)
            
        logging.info("Pipeline finished successfully.")

if __name__ == "__main__":
    generator = VectorGenerator()
    try:
        generator.process_dataset()
    except KeyboardInterrupt:
        logging.info("Process interrupted by user. Progress saved.")
    except Exception as e:
        logging.error(f"Fatal error: {e}")
\end{lstlisting}


\section{Vector Visualization} \label{sec:visualize}
\textbf{Purpose:} We use this script to compute the cosine similarity between the generated vectors and visualize their relationships interactively using \texttt{pyvis}. It filters relationships based on a similarity threshold.

\begin{lstlisting}[style=pythoncode, caption={visualize\_vector\_relationships.py}]
import json
import numpy as np
from sklearn.metrics.pairwise import cosine_similarity
from pyvis.network import Network
import logging

logging.basicConfig(level=logging.INFO, format='%(asctime)s - %(message)s')

VECTOR_FILE = "kurdish_medical_vectors.jsonl"
OUTPUT_HTML = "vector_relationships.html"
SIMILARITY_THRESHOLD = 0.85  # Adjust this to show more/fewer relationships

def load_vectors():
    ids = []
    vectors = []
    logging.info(f"Loading vectors from {VECTOR_FILE}...")
    try:
        with open(VECTOR_FILE, 'r', encoding='utf-8') as f:
            for line in f:
                data = json.loads(line)
                ids.append(str(data['id']))
                vectors.append(data['vector'])
    except FileNotFoundError:
        logging.error(f"{VECTOR_FILE} not found. Please run vector_generator.py first.")
        return [], []
    return ids, np.array(vectors)

def build_graph():
    ids, vectors = load_vectors()
    if not ids:
        return
        
    logging.info(f"Loaded {len(ids)} vectors. Computing similarities...")
    
    # Compute cosine similarity between all vectors
    similarity_matrix = cosine_similarity(vectors)
    
    # Initialize pyvis network
    logging.info("Building interactive graph...")
    net = Network(height="100vh", width="100%", bgcolor="#222222", font_color="white")
    
    # Add nodes
    for i, node_id in enumerate(ids):
        # We only add nodes if they have at least one connection, to keep it clean.
        # But for now, we'll add all of them, or just the ones that connect.
        pass

    # Find relationships (edges) based on similarity using optimized NumPy operations
    edges = []
    connected_nodes = set()
    
    # Get the upper triangle of the similarity matrix (excluding diagonal)
    row_indices, col_indices = np.triu_indices_from(similarity_matrix, k=1)
    
    # Filter indices where similarity is above threshold
    valid_mask = similarity_matrix[row_indices, col_indices] >= SIMILARITY_THRESHOLD
    valid_rows = row_indices[valid_mask]
    valid_cols = col_indices[valid_mask]
    valid_sims = similarity_matrix[valid_rows, valid_cols]
    
    # Optional safety measure: limit max edges so the browser doesn't crash on massive graphs
    MAX_EDGES = 10000
    if len(valid_sims) > MAX_EDGES:
        logging.warning(f"Found {len(valid_sims)} edges, which might freeze the browser. Limiting to top {MAX_EDGES} strongest edges.")
        # Get indices of top MAX_EDGES similarities
        top_indices = np.argsort(valid_sims)[-MAX_EDGES:]
        valid_rows = valid_rows[top_indices]
        valid_cols = valid_cols[top_indices]
        valid_sims = valid_sims[top_indices]

    for r, c, sim in zip(valid_rows, valid_cols, valid_sims):
        edges.append((int(r), int(c), float(sim)))
        connected_nodes.add(int(r))
        connected_nodes.add(int(c))
                
    logging.info(f"Found {len(edges)} strong relationships (similarity >= {SIMILARITY_THRESHOLD}).")
    
    # Add only connected nodes to keep the graph readable
    for node_idx in connected_nodes:
        net.add_node(node_idx, label=f"Document {ids[node_idx]}", title=f"ID: {ids[node_idx]}")
        
    # Add edges
    for i, j, sim in edges:
        # Weight the edge based on similarity
        weight = (sim - SIMILARITY_THRESHOLD) / (1.0 - SIMILARITY_THRESHOLD) * 5
        net.add_edge(i, j, value=weight, title=f"Similarity: {sim:.2f}")

    # Configure physics for a nice Neo4j-like layout
    net.barnes_hut(gravity=-8000, central_gravity=0.3, spring_length=250)
    
    # Save the HTML file
    net.save_graph(OUTPUT_HTML)
    logging.info(f"✅ Graph saved to {OUTPUT_HTML}!")
    logging.info(f"Open {OUTPUT_HTML} in any web browser to view the relationships.")

if __name__ == "__main__":
    build_graph()
\end{lstlisting}


\section{Setup Script} \label{sec:setup}
\textbf{Purpose:} We use this shell script to prepare the environment, install required Python dependencies from \texttt{requirements.txt}, create a \texttt{.env} template, and cache the BERT model.

\begin{lstlisting}[style=bashcode, caption={setup.sh}]
#!/bin/bash

# KMC Vector Generation Setup Script

echo "🚀 Starting setup for KMC Vector Pipeline..."

# 1. Check for Python
if ! command -v python3 &> /dev/null; then
    echo "❌ Python3 is not installed. Please install it first."
    exit 1
fi

# 2. Install Requirements
echo "📥 Installing dependencies from requirements.txt..."
pip install --upgrade pip
pip install -r requirements.txt

# 3. Create .env template if it doesn't exist
if [ ! -f .env ]; then
    echo "📝 Creating .env template..."
    echo "MODEL_NAME=bert-large-cased" > .env
    echo "DATASET_PATH=../just_data_vector/kurdish_medical_corpus_kmc.json" >> .env
    echo "OUTPUT_PATH=kurdish_medical_vectors.jsonl" >> .env
fi

# 4. Download and Cache BERT Model
echo "🧠 Downloading and caching the 'bert-large-cased' model..."
python3 -c "
from transformers import AutoTokenizer, AutoModel
print('Downloading tokenizer...')
AutoTokenizer.from_pretrained('bert-large-cased')
print('Downloading model...')
AutoModel.from_pretrained('bert-large-cased')
print('Model cached successfully!')
"

echo "✅ Setup complete!"
echo "💡 You can now run the pipeline: 'python vector_generator.py'"
\end{lstlisting}


\section{Model Loader} \label{sec:loader}
\textbf{Purpose:} A dedicated script used to download and cache the \texttt{bert-large-cased} model from Hugging Face into the local environment before running the main pipeline.

\begin{lstlisting}[style=bashcode, caption={load\_model\_gpu\_1.sh}]
#!/bin/bash

# Load Model into Cache

echo -e "\033[1;33m[*] Action: Caching 'bert-large-cased' model...\033[0m"

# We use python to trigger the Hugging Face download
python3 -c "
from transformers import AutoTokenizer, AutoModel
print('Downloading tokenizer...')
AutoTokenizer.from_pretrained('bert-large-cased')
print('Downloading model...')
AutoModel.from_pretrained('bert-large-cased')
print('Done!')
"

echo -e "\033[1;32m[SUCCESS] Model 'bert-large-cased' is cached successfully!\033[0m"
\end{lstlisting}

\section{Requirements} \label{sec:requirements}
\textbf{Purpose:} Defines the required Python packages for running the vector generation and visualization pipeline.

\begin{lstlisting}[style=bashcode, caption={requirements.txt}]
transformers
torch
tqdm
scikit-learn
pyvis
\end{lstlisting}

\section{Environment Variables} \label{sec:env}
\textbf{Purpose:} The configuration settings used by the vector generator.

\begin{lstlisting}[style=bashcode, caption={.env}]
MODEL_NAME=bert-large-cased
DATASET_PATH=../just_data_vector/kurdish_medical_corpus_kmc.json
OUTPUT_PATH=kurdish_medical_vectors.jsonl
\end{lstlisting}

\end{document}
